\section*{अध्याय \ref{ch:b3:advanced-nmr} --- NMR गहराई में: स्पंद, विश्रांति और द्विविमीय स्पेक्ट्रम}

\begin{solution}{exo:b3:advanced-nmr:1}
$\nu_0 = (\gamma/2\pi)B_0$: \ce{^1H} \qty{600.4}{MHz}, \ce{^{13}C} \qty{151.0}{MHz}, \ce{^{19}F}
\qty{565.1}{MHz}।
\end{solution}

\begin{solution}{exo:b3:advanced-nmr:2}
$h\nu_0/2kT = 6.626\times10^{-34} \times 400.2\times10^6/(2 \times 1.381\times10^{-23} \times 300)
= \num{3.2e-5}$। \qty{4}{K} पर यह 75 गुना बड़ा है, \num{2.4e-3} (सन्निकटन
अब भी मान्य है)।
\end{solution}

\begin{solution}{exo:b3:advanced-nmr:3}
\qty{10}{\micro s} में $\theta = \gamma B_1t_p = \pi/2$: $\gamma B_1/2\pi = 1/(4 \times
\qty{10}{\micro s}) = \qty{25}{kHz}$। $180^\circ$ स्पंद \qty{20}{\micro s} चलता है; \qty{5}{\micro s}
से $45^\circ$ मिलता है।
\end{solution}

\begin{solution}{exo:b3:advanced-nmr:4}
$T_2 = 1/(\pi \times 3.2) = \qty{0.10}{s}$।
\end{solution}

\begin{solution}{exo:b3:advanced-nmr:5}
$(10.71/42.58)^3 = 0.0159$, गुणा 0.011: \num{1.75e-4}, प्रोटॉन संकेत का लगभग
1/5700। चूँकि संकेत-रव अनुपात $\propto\sqrt N$, समान अनुपात के लिए $5700^2 \approx \num{3e7}$ गुना अधिक
स्कैन चाहिए --- असंभव, इसीलिए ${}^{13}$C स्पेक्ट्रमों में अधिक सांद्र नमूने,
वियुग्मन और ध्रुवण स्थानांतरण प्रयुक्त होते हैं।
\end{solution}

\begin{solution}{exo:b3:advanced-nmr:6}
$T_1 = 1.39/\ln2 = \qty{2.0}{s}$; पुनरावर्तन विलंब $\ge 5T_1 = \qty{10}{s}$।
\end{solution}

\begin{solution}{exo:b3:advanced-nmr:7}
\omterm{def:b3:advanced-nmr:experiments}{COSY}: एक \omterm{def:b3:advanced-nmr:two-d}{तिर्यक शिखर}, एथिल समूह के \ce{CH2} चतुष्क और \ce{CH3} त्रिक के बीच;
\ce{CH3CO} एकक किसी से युग्मित नहीं होता। एकक मेथिल से \omterm{def:b3:advanced-nmr:experiments}{HMBC}:
कार्बोनिल कार्बन (दो आबंध) और \ce{CH2} कार्बन (तीन आबंध)।
\end{solution}

\begin{solution}{exo:b3:advanced-nmr:8}
$k_c = \pi \times 50/\sqrt2 = \qty{111}{s^{-1}}$। $\Delta G^\ddagger = 8.314 \times 330 \times
\ln(6.88\times10^{12}/111) = \qty{68}{kJ/mol}$।
\end{solution}

\begin{solution}{exo:b3:advanced-nmr:9}
NOE $\propto r^{-6}$: $r_b = 0.30 \times 4^{-1/6} = \qty{0.24}{nm}$।
\end{solution}

\begin{solution}{exo:b3:advanced-nmr:10}
अनुनाद पर \omterm{def:b3:advanced-nmr:magnetisation}{घूर्णी निर्देश-तंत्र} में केवल $\mathbf B_1 = B_1\mathbf e_{x'}$ बचता है:
$\dot M_{x'} = 0$, $\dot M_{y'} = \gamma B_1M_z$, $\dot M_z = -\gamma B_1M_{y'}$
($\gamma\mathbf M\times\mathbf B_1$ के घटक)। अतः $M_{x'}$ अचर है और $(M_{y'}, M_z)$ कोणीय
दर $\gamma B_1$ पर घूमता है: $t_p$ के बाद, कोण $\gamma B_1t_p$ से, $x'$ के परितः।
\end{solution}

\begin{solution}{exo:b3:advanced-nmr:11}
थोड़े प्रबल क्षेत्र में स्थित चक्रण
$\tau$ के दौरान एक अचर अतिरिक्त दर से कला प्राप्त करता है; $180^\circ$ स्पंद उसकी कला उलट देता है, और
दूसरे $\tau$ के दौरान वह ठीक उतनी ही खो देता है: ऐसे सभी चक्रण $2\tau$ पर फिर संरेखित हो जाते हैं। यादृच्छिक चक्रण--चक्रण
अन्योन्यक्रियाएँ समय में उच्चावचित होती हैं और दूसरे अंतराल में दोहराई नहीं जातीं: उनका
कला-विचलन निरस्त नहीं होता। प्रतिध्वनि का आयाम वास्तविक $T_2$ के साथ क्षीण होता है।
\end{solution}

\begin{solution}{exo:b3:advanced-nmr:12}
एथिल एथेनोएट, \ce{CH3COOCH2CH3}: \ce{OCH2} चतुष्क (4.12), \ce{CH3CO} एकक (2.05),
\ce{CH3} त्रिक; चतुष्क प्रोटॉन एस्टर ऑक्सीजन के माध्यम से तीन आबंध दूर स्थित
कार्बोनिल कार्बन को देखते हैं। मेथिल प्रोपेनोएट, \ce{CH3CH2COOCH3}, अपना
चतुष्क लगभग 2.3 ppm पर दिखाता (\omterm{def:b3:advanced-nmr:experiments}{HSQC} में लगभग 27 ppm के कार्बन पर, 60 पर नहीं) और एक एकक
\ce{OCH3} लगभग 3.7 ppm पर, जिसका कार्बोनिल तक \omterm{def:b3:advanced-nmr:experiments}{HMBC} शिखर ऑक्सीजन को पार करता है।
\end{solution}

\begin{solution}{pb:b3:advanced-nmr:1}
\textbf{1.} $(2 \times 6 + 2 - 12)/2 = 1$: एक C=O।
\textbf{2.} 173.6 ppm: एस्टर कार्बोनिल; 60.1 ppm: \ce{OCH2}।
\textbf{3.} पाँच; DEPT-135: \ce{CH2} ऋणात्मक (60.1, 36.3, 18.6), \ce{CH3} धनात्मक (14.3,
13.7), कार्बोनिल अनुपस्थित।
\textbf{4.} \qty{400.2}{MHz} और \qty{100.7}{MHz}।
\textbf{5.} $3J = \qty{21.6}{Hz} = \qty{0.054}{ppm}$।
\textbf{6.} बहुकताएँ दिखाती हैं कि हर अणुखंड के भीतर कौन-से समूह पड़ोसी हैं,
परंतु यह नहीं कि अणुखंड ऑक्सीजन और कार्बोनिल के पार कैसे जुड़ते हैं।
\textbf{7.} 4.13--1.25; 2.28--1.66; 1.66--0.95 (और उनके सममित शिखर)।
\textbf{8.} \ce{OCH2CH3} (4.13, 1.25) और \ce{CH2CH2CH3} (2.28, 1.66, 0.95)।
\textbf{9.} \ce{OCH2} और \ce{CH2C=O} प्रोटॉन ऑक्सीजन और
कार्बोनिल कार्बन से अलग हैं: पाँच आबंध, कोई विभेदित युग्मन नहीं।
\textbf{10.} 1.66 ppm, \ce{CH2} और \ce{CH3} के बीच: समान $J$ के साथ 2 + 3 = 5 प्रोटॉनों से
युग्मित, एक षट्क (एक बहुक)।
\textbf{11.} युग्मित मेथिल, अतः सन्निकट कार्बनों पर दो \ce{CH3} (एक \ce{CH3-CH3} इकाई): यहाँ असंभव।
\textbf{12.} \ce{-O-CH2-CH3} और \ce{-CH2-CH2-CH3}, साथ में \ce{C=O}।
\textbf{13.} 4.13/60.1; 2.28/36.3; 1.66/18.6; 1.25/14.3; 0.95/13.7।
\textbf{14.} 14.3 (एथिल \ce{CH3}) और 13.7 ppm (ब्यूटेनॉयल \ce{CH3}), केवल \qty{0.6}{ppm} अलग।
\textbf{15.} 4.13 ppm से 173.6 ppm तक: H--C--O--C, तीन आबंध।
\textbf{16.} 2.28 ppm (दो आबंध) और 1.66 ppm (तीन आबंध) से 173.6 ppm तक।
\textbf{17.} \ce{CH3CH2CH2C(=O)OCH2CH3}, एथिल ब्यूटेनोएट। प्रोपिल प्रोपेनोएट
लगभग 4.0 ppm पर एक \ce{OCH2} त्रिक (चतुष्क नहीं) और लगभग 2.3 ppm पर एक चतुष्क रखता।
\textbf{18.} चार-आबंध युग्मन प्रायः 1 Hz से कम होते हैं; \omterm{def:b3:advanced-nmr:experiments}{HMBC} विलंब, जो
5--10 Hz के लिए समस्वरित है, उन्हें नहीं चुनता।
\textbf{19.} कार्बन भिन्न दरों पर विश्रांति करते हैं और वियुग्मन से भिन्न नाभिकीय
ओवरहाउज़र संवर्धन पाते हैं।
\textbf{20.} $1 - \eu^{-5/20} = 0.22$: कार्बोनिल का प्रतिनिधित्व 4 के गुणक से कम होता है।
\textbf{21.} $1 - \eu^{-5} = 0.993$।
\textbf{22.} प्रोटॉन वियुग्मन से उत्पन्न NOE, जो हर कार्बन के लिए भिन्न है; उसे
प्रतिलोम-द्वारित वियुग्मन (वियुग्मक केवल अधिग्रहण के दौरान चालू) से हटाया जाता है।
\textbf{23.} $256 \times \qty{100}{s} = \qty{25600}{s}$, लगभग 7 घंटे।
\textbf{24.} \textbf{कार्बोनिल कार्बन का $5T_1$: स्कैनों के बीच \qty{100}{s}।}
\end{solution}
